\documentclass[10pt]{article}
\usepackage[margin=22mm]{geometry}
\usepackage{microtype}
\pagestyle{empty}
\setlength{\parskip}{0.45em}
\setlength{\parindent}{0pt}

\begin{document}

\hfill 29 July 2026

Editor\\
\emph{Electric Power Systems Research}

Dear Editor,

Please consider our manuscript, ``GFNO-PINO: A Topology-Disjoint Benchmark
Isolating Physics Fine-Tuning, Topology Conditioning, and Spectral Filtering
in AC-OPF Graph Surrogates,'' for publication in \emph{Electric Power Systems
Research}.

This is a controlled ablation and benchmark study. Its central finding is that
physics-informed fine-tuning substantially improves the power-balance accuracy
of graph-based AC-OPF surrogates under the executed protocol, whereas explicit
topology conditioning and spectral graph filtering---architecturally plausible
choices commonly assumed useful---are not confirmed to improve accuracy in
matched-parameter, topology-disjoint comparisons. We report this negative
result explicitly because it directly informs surrogate-model design choices.

Across five training seeds and four IEEE systems, paired physics fine-tuning
reduces mean sample-maximum power-balance mismatch by 63.4--89.3\%, with all
20 paired seed effects positive. Parameter-matched topology-blind and
spatial-GNN baselines nevertheless match or outperform the full model on three
of four systems.

We do not present the model as a production AC-OPF solver. Every learned model
has 0\% feasibility under the declared proxy, and the manuscript keeps this
limitation prominent. It also distinguishes raw GPU forward-pass latency from
end-to-end feasible-solution time and states that operational use requires
independent verification and feasibility restoration.

Thank you for your consideration.

Sincerely,

\vspace{0.5em}
Roshan Kumar Sharma\\
Ujjwal Prasad Dahal\\
Department of Electrical Engineering\\
Purwanchal Campus, Institute of Engineering\\
Tribhuvan University, Dharan, Nepal

\end{document}
